% dependency-management.tex — package managers for iOS + RN (SPM, CocoaPods incl. trunk going read-only,
% Carthage, npm / Yarn / pnpm), semver and range operators, lockfiles, transitive + peer dependencies (a vulnerable
% transitive dependency drawn with the fix ladder), choosing a library (maintenance signals, bus factor, size,
% licence compatibility), vendoring vs depending, update strategy with Renovate / Dependabot, removing deps.
% Attack side (npm incidents, install scripts, provenance, SBOM, cooldowns) lives in security/supply-chain-security.tex;
% the weighted build-vs-buy scorecard in engineering/tech-debt-and-decision-records.tex — only pointed to.
% Sources (checked 2026-10-06) — CocoaPods trunk date re-read; the rest are specs/docs:
% https://semver.org/ · https://github.com/npm/node-semver#advanced-range-syntax (caret / tilde, 0.x rules)
% https://docs.npmjs.com/cli/v10/configuring-npm/package-json#overrides · https://docs.npmjs.com/cli/v10/commands/npm-explain
% https://docs.swift.org/package-manager/ (from: = up to next major, .upToNextMinor, Package.resolved)
% https://blog.cocoapods.org/CocoaPods-Specs-Repo/ (trunk read-only from 2 Dec 2026; existing builds keep working)
% https://pnpm.io/motivation (content-addressable store, strict node_modules)
% https://www.apache.org/licenses/LICENSE-2.0 (patent grant, NOTICE) · https://www.gnu.org/licenses/license-list.html
% https://docs.renovatebot.com/ · https://docs.github.com/en/code-security/dependabot
% @labels: area=engineering kind=tooling platform=cross-platform topic=build,tooling discussion=128
% @tags: swift-package-manager, cocoapods, npm, pnpm, semver, caret-and-tilde-ranges, lockfile, transitive-dependencies, peer-dependencies, npm-overrides, open-source-licences, renovate, dependabot, vendoring
\documentclass[8pt]{extarticle}
\usepackage{printup-sheet}
\usepackage{array}

\lstdefinelanguage{DepSheet}{
  morekeywords={dependencies,overrides,package,from,upToNextMinor,exact},
  sensitive=true, morecomment=[l]{//}, morestring=[b]",
  literate={->}{{\hbox{-}\hbox{>}}}2 {>=}{{\hbox{>}\hbox{=}}}2 {<=}{{\hbox{<}\hbox{=}}}2
           {..<}{{\hbox{.}\hbox{.}\hbox{<}}}3 {==}{{\hbox{=}\hbox{=}}}2}
\lstset{basicstyle=\ttfamily\scriptsize, aboveskip=2pt, belowskip=2pt}
\newcommand\ct[1]{\texttt{#1}}
\newcommand\hd[1]{\par\vspace{3pt}\noindent{\bfseries\color{sheetBlue}#1}\par\vspace{1pt}}
\newcommand\gee{\hbox{>}\hbox{=}}
\newcommand\hrr{\hbox{.}\hbox{.}\hbox{<}}
\newcommand\crt{\textasciicircum}
\newcommand\tld{\textasciitilde}

\tikzset{
  pk/.style={box, font=\scriptsize, inner sep=2pt, minimum width=17mm, align=center},
  bad/.style={pk, draw=sheetRed, fill=sheetRed!10, very thick},
  l/.style={font=\tiny, text=black!75, inner sep=1pt, align=center},
  st/.style={font=\tiny, inner sep=2pt, align=left, text width=38mm, draw=#1, fill=#1!6, rounded corners=1pt, anchor=west},
}

\begin{document}

\sheettitle{Dependency management — versions, lockfiles, choosing, updating}{engineering · folio}

\oneliner{Every dependency is code you ship, cannot fully read, and must keep updating. \textbf{Semver ranges}
say what you \emph{accept}; the \textbf{lockfile} records what you \emph{got} (every transitive package, exact
version + hash) so every machine and CI builds the same tree. Choose libraries by \textbf{maintenance, size,
licence and exit cost}, update in \textbf{small automated steps}, and remember that most of your tree is
\textbf{transitive} — you are responsible for packages you never chose.}

\vspace{3pt}
\noindent\begin{tikzpicture}[sheet]
  \node[font=\bfseries\small, text=sheetBlue, anchor=west] at (-0.1,3.55) {A vulnerable transitive dependency — and the fix ladder};
  \node[pk, draw=sheetBlue, fill=sheetBlue!15] (app) at (3.7,2.95) {\textbf{App}\\\ct{package.json}};
  \node[pk] (a) at (1.2,1.85) {\ct{A} \ct{\crt2.1.0}\\$\to$ 2.3.0};
  \node[pk] (b) at (3.7,1.85) {\ct{B} \ct{\tld1.4.0}\\$\to$ 1.4.7};
  \node[pk] (e) at (6.2,1.85) {\ct{E} 3.0.1\\(exact)};
  \node[bad] (c) at (2.45,0.55) {\ct{C} \textbf{1.4.2}\\CVE: fixed in 1.4.3};
  \node[pk] (d) at (5.0,0.55) {\ct{D} 0.9.1};
  \draw[flow] (app) -- node[l, left]{direct} (a); \draw[flow] (app) -- (b); \draw[flow] (app) -- (e);
  \draw[flow] (a) -- node[l, left]{\ct{\crt1.4.0}} (c); \draw[flow] (b) -- node[l, right]{\ct{1.4.2} exact} (c);
  \draw[flow] (b) -- (d); \draw[flow] (e) -- (d);
  \node[l, text=sheetRed] at (2.45,-0.05) {transitive: you never chose it, you ship it};
  \node[l, text=sheetGrey, align=left, anchor=west] at (6.9,0.55) {lockfile pins\\all 5, with\\integrity hashes};
  \draw[sheetGrey!40] (8.35,3.7) -- (8.35,-0.2);
  % fix ladder
  \node[st=sheetBlue] at (8.55,3.05) {\textbf{1 find the path}: \ct{npm explain C} / \ct{npm ls C} · \ct{swift package show-dependencies}};
  \node[st=sheetGreen!70!black] at (8.55,2.35) {\textbf{2 update the direct dep}: a B release that allows C $\geq$ 1.4.3 — best: upstream fixed};
  \node[st=sheetOrange] at (8.55,1.65) {\textbf{3 override} the transitive: npm \ct{overrides}, Yarn \ct{resolutions}, pnpm \ct{overrides}; test, add a removal note};
  \node[st=sheetBrown] at (8.55,0.95) {\textbf{4 patch / fork} (\ct{patch-package}, a pinned fork) — you now own it};
  \node[st=sheetRed] at (8.55,0.3) {\textbf{5 remove} the feature or the dependency};
  \node[l, text=sheetGrey, align=left, anchor=west] at (12.75,2.7) {not exploitable in your\\usage? record why, with\\an expiry — don't just\\silence the scanner};
  \node[l, text=sheetBrown, align=left, anchor=west] at (12.75,1.3) {SPM: \emph{one} version of C per\\graph — incompatible ranges\\fail resolution. npm: may\\nest two copies side by side};
\end{tikzpicture}

\begin{multicols}{2}
\raggedright\setstretch{1.0}

\hd{Package managers}
{\footnotesize\setlength{\tabcolsep}{3pt}
\begin{tabular}{@{}>{\raggedright\arraybackslash}p{13mm}>{\raggedright\arraybackslash}p{20mm}>{\raggedright\arraybackslash}p{41mm}@{}}
\toprule
\textbf{Tool} & \textbf{Lockfile} & \textbf{Model} \\ \midrule
SPM & \ct{Package.resolved} & built into Xcode/\ct{swift}; source + \ct{binaryTarget} (XCFramework + checksum); one version per package \\
CocoaPods & \ct{Podfile.lock} & Ruby; generates a workspace; trunk \textbf{read-only from 2 Dec 2026} (existing pods still install) — RN still uses it \\
Carthage & \ct{Cartfile.resolved} & builds frameworks, you link them; decentralised, rarely chosen now \\
npm & \ct{package-lock.json} & nested \ct{node\_modules}: two versions can coexist \\
Yarn & \ct{yarn.lock} & Berry adds Plug'n'Play; RN projects use \ct{nodeLinker: node-modules} \\
pnpm & \ct{pnpm-lock.yaml} & content-addressed store + symlinks; \textbf{strict}: undeclared (``phantom'') imports fail \\
\bottomrule
\end{tabular}}

\hd{Semver and ranges}
\begin{itemize}
  \item \ct{MAJOR.MINOR.PATCH} $=$ breaking . feature . fix; \ct{0.y.z}: anything may change. A promise by
        humans, not a guarantee — tests decide.
  \item npm: \ct{\crt1.2.3} $=$ \ct{\gee1.2.3 <2.0.0} · \ct{\crt0.2.3} $=$ \ct{\gee0.2.3 <0.3.0} (caret treats the first
        non-zero as major) · \ct{\tld1.2.3} $=$ \ct{\gee1.2.3 <1.3.0} · \ct{1.2.3} exact. Pre-releases
        (\ct{2.0.0-beta.1}) are skipped by ranges unless asked for.
  \item SPM: \ct{from: "1.2.3"} $=$ \ct{1.2.3\hrr2.0.0} · \ct{.upToNextMinor(from:)} $=$ \ct{1.2.3\hrr1.3.0} ·
        \ct{exact:} · branch/revision (not for releases).
  \item \textbf{Apps}: pin native-heavy deps exactly (RN, SDKs) and let the lockfile pin the rest.
        \textbf{Libraries}: wide ranges, peers for shared singletons (\ct{react}, \ct{react-native}).
\end{itemize}

\hd{Lockfile rules}
Commit it (apps); CI installs \emph{from} it (\ct{npm ci}, \ct{pod install}, resolved SPM) and fails if it
drifts. Review lockfile diffs: a one-line change in \ct{package.json} can move 200 transitive versions.
A library's lockfile is not published — consumers resolve its ranges themselves.

\hd{Example — ranges, override, SPM}
\begin{lstlisting}[language=DepSheet]
// package.json (comments for the sheet only)
"dependencies": {
  "react-native": "0.81.0",      // exact: native, upgrade on purpose
  "date-fns": "^4.1.0",          // >=4.1.0 <5.0.0
  "@example/analytics": "~2.3.1" // >=2.3.1 <2.4.0
},
"overrides": { "semver": "7.5.4" } // force a patched transitive
// Package.swift
.package(url: "https://example.com/a.git", from: "1.1.0")
                                        // 1.1.0..<2.0.0
.package(url: "https://example.com/b.git",
         .upToNextMinor(from: "0.9.2")) // 0.9.2..<0.10.0
\end{lstlisting}

\hd{Remember}
\textbf{``Ranges say what you accept, the lockfile what you got; every dependency is yours to update — or remove.''}

\columnbreak

\hd{Choosing a library — the signals}
\begin{itemize}
  \item \textbf{Maintenance}: releases in the last months, issues and PRs answered, CI green, a changelog,
        honest semver; supports the current OS, Swift 6 / RN New Architecture, ships types.
  \item \textbf{Bus factor}: more than one active maintainer, or a company behind it; who fixes a CVE at 2 am?
  \item \textbf{Size + cost}: binary size delta (App Thinning Size Report, JS bundle analyser), transitive
        count, launch cost (dynamic frameworks), permissions and privacy manifest it brings.
  \item \textbf{Need it?} 20 lines of your own beat a package for a helper; platform APIs (\ct{URLSession}
        async/await, \ct{fetch}, \ct{Intl}) replace many classics. \textbf{Exit cost}: wrap it behind your interface.
\end{itemize}

\hd{Licences — what they cost an app}
{\footnotesize\setlength{\tabcolsep}{3pt}
\begin{tabular}{@{}>{\raggedright\arraybackslash}p{19mm}>{\raggedright\arraybackslash}p{55mm}@{}}
\toprule
MIT, BSD & permissive: keep copyright + licence text (Acknowledgements screen) \\
Apache-2.0 & permissive + \textbf{patent grant}; keep \ct{NOTICE}; not GPLv2-compatible \\
MPL-2.0 & file-level copyleft: publish changes to \emph{its} files only \\
LGPL & users must be able to relink a modified library — hard with static linking in a signed iOS app:
       legal review \\
GPL / AGPL & the combined app must ship as GPL source (AGPL: network use too) — a no for closed apps \\
no licence & all rights reserved: you may not use it \\
\bottomrule
\end{tabular}}\par
Gate it in CI: a licence scanner emits SPDX ids, checked against an allow-list.

\hd{Vendor, update, remove}
\begin{itemize}
  \item \textbf{Vendor} (copy source/binary into the repo) for tiny, abandoned or patched code, or
        offline builds — you now own its updates and CVEs. Middle ground: pinned fork, \ct{patch-package}.
  \item \textbf{Update}: Renovate / Dependabot PRs on a schedule; group patch + minor, majors one by one with
        the changelog read; automerge only low-risk groups on green CI. RN/Xcode upgrades are planned work.
  \item \textbf{Remove}: find imports, swap in the platform API, delete, check the lockfile and binary
        shrank; \ct{knip} / \ct{depcheck} find unused JS deps.
\end{itemize}

\hd{Traps}
\begin{itemize}
  \trap{Two copies of \ct{react} in an RN monorepo $\to$ ``Invalid hook call'': dedupe; libraries declare it as a \textbf{peer}.}
  \trap{``We pinned \ct{\crt1.2.3}'' — a caret range is not a pin; the lockfile is.}
  \trap{Code relying on a transitive package it never declared — breaks under pnpm or the next update.}
\end{itemize}


\end{multicols}

\noindent{\footnotesize\color{sheetGrey}\textit{Related:} supply-chain-security (incidents, provenance, SBOM) ·
tech-debt-and-decision-records (scorecard) · modularization-spm · rn-tooling-release · large-migrations}

\end{document}
